\section{Conclusions}
\begin{frame}{Conclusions}
\framesubtitle{Expected Contributions}
\begin{itemize}
    \item This work motivates the use of decentralized certification protocols as an alternative to centralized document verification systems.
    \item The literature analysis found that Bitcoin-based approaches face limitations in scalability, privacy, and applicability to document verification.
    \item Existing solutions based on \texttt{OP\_RETURN} and OpenTimestamps exhibit structural limitations at the protocol level.
    \item This work proposes a Taproot-based approach using Bitcoin-native cryptographic tools: Schnorr and Mast.
\end{itemize}
\end{frame}

\begin{frame}{Conclusions}
\framesubtitle{Expected Technical Capabilities}
\begin{itemize}
    \item The proposed protocol is designed to support integrity verification through hash commitments anchored in Merkle trees.
    \item Issuer authenticity is expected to be verified using crytographic digital signatures.
    \item Document revocation is addressed through transaction chaining and bitmaps on op return.
    \item Privacy is preserved by batching documents without revealing any information about unrelated documents.
\end{itemize}
\end{frame}


\begin{frame}{Conclusions}
\framesubtitle{Limitations and Future Work}
\begin{itemize}
    \item The proposed protocol does not address the semantic validity or correctness of document content.
    \item Document revocation remains under the control and responsibility of the issuing institution.
    \item Long-term document preservation and proof availability after issuance are not considered.
    \item The study of long-term preservation mechanisms remains an open area for future work.
    \item The design of incentives or penaltys mechanisms for misuse are left for future research.
\end{itemize}
\end{frame}
